\begin{abstract}
Official crop yield statistics are often reported at the level of administrative units, such as municipalities in Brazil and counties in the United States, rather than at finer spatial resolutions.
However, precision agriculture, insurance, and localized risk assessment require estimates at finer spatial resolutions.
We propose MSTNet (\emph{Multi} STNet), a Transformer model that combines data sampled at different temporal frequencies.
Separate STNet trunks encode soy-masked Sentinel-2 time series and a daily gridded climate series, while annual soil properties enter the decoder by late fusion.
Since the model is able to predict a yield value per soy pixel, we can use it to predict yield for arbitrary polygons by averaging the pixel-wise predictions.
Training uses the soy-masked municipality-mean Sentinel-2 and climate series, as well as the annual soil features under PAM supervision;
we also evaluate the prediction of every soy pixel of a given municipality and average it, to evaluate how well the model can generalize to arbitrary polygons.
We evaluate in Paran\'{a}, Brazil, with year leave-one-out over harvest years 2020--2024, reporting only the left-out year in each fold.
Excluding the 2022 drought fold, the pixel-wise average LOO $R^2$ is $0.302$; when 2022 itself is the left out year, performance collapses because no comparable drought year remains in training dataset.
Direct comparison with earlier work is limited for this municipal-supervision, pixel-scoring setting; we report the LOO benchmark and leave field-referenced checks for future work.
\end{abstract}
